\SourcePage{1}
\begin{center}
{\Large\sffamily\bfseries \term{perverse-sheaf}{反常层}}\par\medskip
{\large Faisceaux pervers}\par\bigskip
Alexander Beilinson\quad Joseph Bernstein\par
Pierre Deligne\quad Ofer Gabber\par\bigskip
Astérisque，第100卷\par
法国数学会，2018年\par
在法国国家科学研究中心资助下出版
\end{center}
\vfill
\noindent\textbf{译本说明}\par
本译本以所提供的2018年再版PDF为底本，完整收录1982年原文及2018年新增材料。原文的条目、公式及参考文献编号予以保留；本译本重新分页。勘误中的页码、行号均指原文，所指版本依勘误的说明确定。正文技术名词与书末中—英—法对照表之间设置双向链接。
\clearpage
\SourcePage{2}
\noindent\textbf{作者与通讯地址}\par
\noindent A. Beilinson：美国芝加哥大学数学系，5734 S. University Avenue, Chicago, IL 60637, USA。\par
电子邮件：\href{mailto:sasha@math.uchicago.edu}{sasha@math.uchicago.edu}
\par\medskip\noindent J. Bernstein：以色列特拉维夫大学数学科学学院，Tel Aviv 69978, Israel。\par
电子邮件：\href{mailto:bernstei@post.tau.ac.il}{bernstei@post.tau.ac.il}
\par\medskip\noindent P. Deligne：美国高等研究院数学学院，1 Einstein Drive, Princeton, NJ 08540, USA。\par
电子邮件：\href{mailto:deligne@ias.edu}{deligne@ias.edu}
\par\medskip\noindent O. Gabber：法国高等科学研究所，Le Bois-Marie, 35 route de Chartres, 91440 Bures-sur-Yvette, France。\par
电子邮件：\href{mailto:gabber@ihes.fr}{gabber@ihes.fr}
\par\bigskip\noindent\textbf{数学主题分类（2010）}
\begin{description}[leftmargin=4.4em,labelwidth=4em]
\item[{18E30}] \term{derived-category}{导出范畴}，\term{triangulated-category}{三角范畴}；
\item[{19F27}] \term{etale}{平展}\term{cohomology}{上同调}；
\item[{32S60}] \term{stratification}{分层}，\term{constructible}{可构造}\term{sheaf}{层}，\term{intersection-cohomology}{交叉上同调}；
\item[{55N33}] \term{intersection-homology}{交叉同调}与\term{intersection-cohomology}{交叉上同调}；
\item[{32S30}] 奇点的形变，\term{vanishing-cycles}{消失轮换}。
\end{description}
\noindent\textbf{关键词}\quad \term{triangulated-category}{三角范畴}，\term{derived-category}{导出范畴}，\term{t-structure}{$t$-结构}，\term{t-structure}{$t$-结构}的\term{heart}{心}，\term{stratification}{分层}，\term{constructible}{可构造}\term{sheaf}{层}，\term{l-adic}{$\ell$-进}\term{sheaf}{层}，\term{perverse-sheaf}{反常层}，\term{intermediate-extension}{中间延拓}函子，\term{self-dual}{自对偶}\term{perversity}{反常度}，\term{vanishing-cycles}{消失轮换}，\term{intermediate-extension}{中间延拓}的\term{pure}{纯}性定理，\term{decomposition-theorem}{分解定理}，相对\term{hard-lefschetz}{强Lefschetz定理}，\term{weight}{权}，\term{pure}{纯}\term{complex}{复形}，\term{mixed}{混合}\term{complex}{复形}。
\par\bigskip
法国数学会感谢Mathdoc团队通过NUMDAM计划为本版出版提供帮助。
\UnnumberedChapter{内容提要}
\SourcePage{3}
\noindent\textbf{法文提要译文}\par
本卷介绍\term{perverse-sheaf}{反常层}理论。第1章给出\term{triangulated-category}{三角范畴}上\term{t-structure}{$t$-结构}的定义及基本性质。第2章分别在\term{stratification}{分层}空间和\term{scheme}{概形}的框架下，引入任意\term{perversity}{反常度}的\term{perverse-sheaf}{反常层}及\term{intermediate-extension}{中间延拓}函子。第3章讨论若干补充内容，包括\term{filtered-derived-category}{滤过导出范畴}与\term{realization-functor}{实现函子}，以及\term{sheaf}{层}的\term{derived-category}{导出范畴}中的\term{localization}{局部化}。第4章汇集\term{self-dual}{自对偶}\term{perversity}{反常度}下\term{perverse-sheaf}{反常层}的基本性质。第5章是全书的核心，研究有限域上簇的\term{mixed}{混合}\term{l-adic}{$\ell$-进}\term{perverse-sheaf}{反常层}，特别证明\term{intermediate-extension}{中间延拓}的\term{pure}{纯}性定理、\term{decomposition-theorem}{分解定理}和相对\term{hard-lefschetz}{强Lefschetz定理}。第6章说明如何将前一章的结果用于复代数几何。本版另收录勘误与增补、新增参考文献，以及讨论若干常用函子\term{t-exact}{$t$-正合}性的附录。
\par\bigskip\noindent\textbf{英文提要译文}\par
本卷阐述\term{perverse-sheaf}{反常层}理论。第1章给出\term{triangulated-category}{三角范畴}上\term{t-structure}{$t$-结构}的定义和基本性质。第2章在\term{stratification}{分层}空间及\term{scheme}{概形}两种背\term{site}{景}下，介绍任意\term{perversity}{反常度}的\term{perverse-sheaf}{反常层}和\term{intermediate-extension}{中间延拓}函子。第3章补充讨论\term{filtered-derived-category}{滤过导出范畴}、\term{realization-functor}{实现函子}，以及\term{sheaf}{层}的\term{derived-category}{导出范畴}中的\term{localization}{局部化}。第4章汇集关于中间\term{perversity}{反常度}的\term{perverse-sheaf}{反常层}的基本事实。作为全书核心的第5章研究有限域上簇的\term{mixed}{混合}\term{l-adic}{$\ell$-进}\term{perverse-sheaf}{反常层}，尤其包含\term{intermediate-extension}{中间延拓}的\term{pure}{纯}性定理、\term{decomposition-theorem}{分解定理}和相对\term{hard-lefschetz}{强Lefschetz定理}。第6章说明如何在复代数几何中使用第5章的结果。本版附有勘误与增补、新增参考文献，以及关于若干常用函子\term{t-exact}{$t$-正合}性的附录。
\SourcePage{4}
\UnnumberedChapter{2018年再版目录}
\SourcePage{5}
\noindent 下列页码沿用再版所印目录；正文中的原印刷页码另见1982年原版目录。
\begin{longtable}{@{}p{.84\textwidth}r@{}}
\toprule 第二版序言 & 1\\
0. 引言 &3\\
1. \term{triangulated-category}{三角范畴}的\term{abelian-subcategory}{阿贝尔子范畴} &14\\
\quad 1.1 \term{triangulated-category}{三角范畴} &14\\
\quad 1.2 \term{abelian-subcategory}{阿贝尔子范畴} &23\\
\quad 1.3 $t$-范畴 &25\\
\quad 1.4 \term{recollement}{粘合} &39\\
2. \term{stratification}{分层}空间与\term{scheme}{概形}上的\term{perverse-sheaf}{反常层} &52\\
\quad 2.1 \term{stratification}{分层}空间 &52\\
\quad 2.2 \term{scheme}{概形} &62\\
3. 补充材料 &72\\
\quad 3.1 \term{filtered-derived-category}{滤过导出范畴}、\term{canonical-filtration}{典范滤过}与\term{stupid-filtration}{朴素滤过} &72\\
\quad 3.2 \term{localization}{局部化} &81\\
\quad 3.3 整系数\term{cohomology}{上同调} &94\\
4. \term{self-dual}{自对偶}\term{perversity}{反常度}：几何性质 &97\\
\quad 4.1 \term{affine}{仿射}态射 &98\\
\quad 4.2 正合性与\term{adjunction}{伴随} &103\\
\quad 4.3 \term{simple-object}{单对象} &108\\
\quad 4.4 \term{vanishing-cycles}{消失轮换}（上界估计） &110\\
\quad 4.5 \term{betti-number}{Betti数}的估计 &113\\
5. \term{self-dual}{自对偶}\term{perversity}{反常度}：\term{weight}{权} &118\\
\quad 5.1 回顾[1] &118\\
\quad 5.2 一个逆命题 &125\\
\quad 5.3 \term{weight-filtration}{权滤过}（原目录误印为4.3） &130\\
\quad 5.4 \term{pure}{纯}\term{complex}{复形}（原目录误印为4.4） &136\\
\SourcePage{6}6. 从有限域到复数域 &144\\
\quad 6.1 原理 &144\\
\quad 6.2 例子 &155\\
参考文献 &162\\
术语索引 &164\\
符号索引 &166\\
附录 &167\\
勘误与增补 &173\\
新增参考文献 &179\\\bottomrule
\end{longtable}
\UnnumberedChapter{第二版序言}
\SourcePage{7}
本卷是《\term{perverse-sheaf}{反常层}》的再版。该书最初作为1981年7月6日至10日在Luminy举行的“奇异空间上的分析与拓扑”会议论文集的第一卷，收入Astérisque第100卷。本版保留原文不变，另补入一篇简短附录、一份勘误与增补\footnote{原版页边的每个数字，均指向原版第173页起勘误与增补中的相应条目。译本将这些材料完整收入书后；其中的定位仍依原版。}，以及新增参考文献。

如引言开头所述，Ofer Gabber当年不愿以Astérisque第100卷合著者的身份署名。此次他加入再版的作者行列，我们深感欣喜。
